\chapter{Grundlagen und Stand der Technik}
\label{chap:background}

\section{Stand der Forschung}
\offen{Relevante Originalpublikationen vergleichen und die daraus
abgeleitete Forschungslücke belegen.}
Hier wird nur ein Zitierbeispiel gesetzt; der nachfolgende Beleg ist ein
ausdrücklich zu ersetzender Quellenplatzhalter.\mcicite{quelle-ersetzen}
Quellen werden möglichst unmittelbar nach der übernommenen Aussage
angegeben und im Verzeichnis nach ihrer Erstnennung geordnet.

\section{Modell und Formelzeichen}
Die folgenden Kapitel enthalten ein eigenständig gekennzeichnetes
Vorlagenbeispiel mit synthetischen Daten. Es werden keine durchgeführten
Messungen behauptet. Als Satzbeispiel wird in \eqnref{eq:ohm} ein ohmsches
Modell bei konstanter Temperatur verwendet:
\begin{equation}
  R = \frac{U}{I}.
  \label{eq:ohm}
\end{equation}
Dabei steht $R$ für den elektrischen Widerstand in \unit{\ohm}, $U$ für
die Spannung in \unit{\volt} und $I$ für die Stromstärke in \unit{\ampere}.
\offen{Das für die Forschungsfrage geeignete Modell fachlich belegen,
seine Annahmen prüfen und alle verwendeten Zeichen erläutern.}

Kurze Direktzitate können mit \mciquote{OFFEN: exakter Originalwortlaut}
gesetzt werden. Quelle und genaue Fundstelle werden ergänzt; auch
Paraphrasen und übernommene Gliederungen erhalten einen Beleg.
